% Zone „Das kennst du schon“ – aus bank/punkte-und-strecken-im-koordinatensystem/zone.jsonl (zusammenbau v0.1)
\einheitenkopf[zone][Kennst du schon]{Das kennst du schon}
\setcounter{aufgabe}{0}

%% TODO Titel: Ich-kann-Satz fehlt in der Bank (Platzhalter: Kettenname) – Nr. 1
%% TODO Anweisung: ein Satz über der ersten Teilaufgabe fehlt in der Bank – Nr. 1
\begin{aufgabe}{Punkte im ebenen Koordinatensystem mit vier Quadranten lesen und eintragen, Vorzeichen der Koordinaten deuten}
\begin{teile}
\teil Trage $A(3 | 2)$ ein. Lies $B$ ab.

\begin{ksys}[xmin=-5,xmax=5,ymin=-5,ymax=5] \punkt{1}{4}{B} \end{ksys}
B( \leerfeld | \leerfeld )
\teil Trage $E(-3 | 2)$ und $F(1{,}5 | -2)$ ein. Lies $G$ ab.

\begin{ksys}[xmin=-5,xmax=5,ymin=-5,ymax=5] \punkt{-2}{-1}{G} \end{ksys}
G( \leerfeld | \leerfeld )
\end{teile}
\end{aufgabe}

%% TODO Titel: Ich-kann-Satz fehlt in der Bank (Platzhalter: Kettenname) – Nr. 2
%% TODO Anweisung: ein Satz über der ersten Teilaufgabe fehlt in der Bank – Nr. 2
\begin{aufgabe}{Schrägbilder von Quadern, Prismen und Pyramiden lesen und skizzieren, Netze aufklappen}
\begin{teile}
\teil Quader im Bild: Wie viele Kanten, wie viele davon verdeckt?

\quader{4}{2}{3}{a}{b}{c}
Kanten: \leerfeld , verdeckt: \leerfeld
\teil Zeichne ein Schrägbild des Quaders mit $5$ cm Breite, $3$ cm Tiefe und $2$ cm Höhe (Tiefe schräg und halb so lang).

\begin{ksys}[xmin=0,xmax=8,ymin=0,ymax=5] \end{ksys}
\end{teile}
\end{aufgabe}

%% TODO Titel: Ich-kann-Satz fehlt in der Bank (Platzhalter: Kettenname) – Nr. 3
%% TODO Anweisung: ein Satz über der ersten Teilaufgabe fehlt in der Bank – Nr. 3
\begin{aufgabe}{Verbindungsvektor bilden (Spitze minus Fuß), Vektoren addieren und vervielfachen, den Betrag berechnen}
\begin{teile}
\teil $A(2 | 1 | 3)$, $B(5 | 5 | 3)$: $\overrightarrow{AB}$ und seine Länge?
\teil $\vec a = (2 | -1 | 3)$, $\vec b = (-4 | 0 | 1{,}5)$: $2\vec a + \vec b$?
\end{teile}
\end{aufgabe}

%% TODO Titel: Ich-kann-Satz fehlt in der Bank (Platzhalter: Kettenname) – Nr. 4
%% TODO Anweisung: ein Satz über der ersten Teilaufgabe fehlt in der Bank – Nr. 4
\begin{aufgabe}{Satz des Pythagoras anwenden und mit Wurzeln rechnen (Wurzel isolieren, dann quadrieren)}
\begin{teile}
\teil Katheten $6$ cm und $8$ cm: Hypotenuse? \leerfeld[cm]
\teil Hypotenuse $2{,}5$ m, eine Kathete $1{,}5$ m: andere Kathete? \leerfeld[m]
\end{teile}
\end{aufgabe}

%% TODO Titel: Ich-kann-Satz fehlt in der Bank (Platzhalter: Kettenname) – Nr. 5
%% TODO Anweisung: ein Satz über der ersten Teilaufgabe fehlt in der Bank – Nr. 5
\begin{aufgabe}{Flächenformeln (Dreieck, Trapez), Prozentrechnung und Strahlensätze für die Sachanhänge}
\begin{teile}
\teil Dreieck mit Grundseite $6$ cm und Höhe $4$ cm: Flächeninhalt? \leerfeld[cm²]
\teil Trapez mit $a = 5{,}5$ m, $c = 2{,}5$ m, $h = 1{,}2$ m: Flächeninhalt? \leerfeld[m²]
\end{teile}
\end{aufgabe}

%% TODO Titel: Ich-kann-Satz fehlt in der Bank (Platzhalter: Kettenname) – Nr. 6
%% TODO Anweisung: ein Satz über der ersten Teilaufgabe fehlt in der Bank – Nr. 6
\begin{aufgabe}{Skalarprodukt in Koordinaten berechnen und Orthogonalität als „Skalarprodukt null“ erkennen}
\begin{teile}
\teil $(2 | 1 | 3) \circ (1 | 4 | 2)$? \leerfeld
\teil $(2 | -3 | 1) \circ (-1 | 2 | 0{,}5)$? \leerfeld
\end{teile}
\end{aufgabe}

%% TODO Titel: Ich-kann-Satz fehlt in der Bank (Platzhalter: Kettenname) – Nr. 7
%% TODO Anweisung: ein Satz über der ersten Teilaufgabe fehlt in der Bank – Nr. 7
\begin{aufgabe}{Eigenschaften der Dreiecke und Vierecke kennen (gleichschenklig, gleichseitig, rechtwinklig; Parallelogramm, Rechteck, Raute, Quadrat, Trapez, Drachenviereck) und den Satz des Thales}
\begin{teile}
\teil Welches Viereck hat vier gleich lange Seiten und vier rechte Winkel? \\ \kreuz{Raute} \\ \kreuz{Quadrat} \\ \kreuz{Rechteck}
\teil Zwei Paare paralleler Gegenseiten, kein rechter Winkel, verschieden lange Nachbarseiten: welches Viereck? \\ \kreuz{Trapez} \\ \kreuz{Parallelogramm} \\ \kreuz{Drachenviereck}
\end{teile}
\end{aufgabe}

%% TODO Titel: Ich-kann-Satz fehlt in der Bank (Platzhalter: Kettenname) – Nr. 8
%% TODO Anweisung: ein Satz über der ersten Teilaufgabe fehlt in der Bank – Nr. 8
\begin{aufgabe}{Kollinearität zweier Vektoren prüfen (Vielfaches)}
\begin{teile}
\teil Ist $(4 | 2 | 6)$ ein Vielfaches von $(2 | 1 | 3)$? Faktor? Faktor \leerfeld
\teil Ist $(4 | -2 | -1)$ ein Vielfaches von $(-2 | 1 | 0{,}5)$? Faktor? Faktor \leerfeld
\end{teile}
\end{aufgabe}

%% TODO Titel: Ich-kann-Satz fehlt in der Bank (Platzhalter: Kettenname) – Nr. 9
%% TODO Anweisung: ein Satz über der ersten Teilaufgabe fehlt in der Bank – Nr. 9
\begin{aufgabe}{Verbindungsvektor bilden (Spitze minus Fuß), Vektoren addieren und vervielfachen, den Betrag berechnen – Fehler finden}
\begin{teile}
\teil Ich finde den Fehler: $U(1 | 2 | 2)$, $V(3 | 3 | 4)$. Ali rechnet $|\overrightarrow{UV}| = \sqrt{3^2 + 3^2 + 4^2} = \sqrt{34}$.
\end{teile}
\end{aufgabe}

%% TODO Titel: Ich-kann-Satz fehlt in der Bank (Platzhalter: Kettenname) – Nr. 10
%% TODO Anweisung: ein Satz über der ersten Teilaufgabe fehlt in der Bank – Nr. 10
\begin{aufgabe}{Verbindungsvektor bilden (Spitze minus Fuß), Vektoren addieren und vervielfachen, den Betrag berechnen – selbst rechnen}
\begin{teile}
\teil $K(2 | 1 | 5)$, $L(6 | 9 | 6)$: Länge von $KL$? \leerfeld
\end{teile}
\end{aufgabe}
